\documentclass[11pt]{article}
\usepackage{fontspec,amsmath,amssymb,amsthm,hyperref}
\usepackage[margin=1in]{geometry}
\setlength{\emergencystretch}{3em}
\title{Relative scores turn the weak equation into entropy production}
\author{LangevinLean mathematical exposition}
\date{}
\begin{document}\maketitle
\noindent Source module: \texttt{UnadjustedLangevin/LSI/ScoreCalculus.lean}. Compile this proof document with XeLaTeX or LuaLaTeX. The accompanying HTML pairs the mathematical argument with the exact Lean statements and records the scope of the formalization.


The chain rule \(\nabla\log r=\nabla r/r\) converts the entropy test \(\log r+1\) into the relative score. The weak identity supplied for this test then gives Fisher dissipation, with an extra score/error pairing when drift forcing is present.


\section{gradient log of pos}


The logarithm chain rule gives \(D\log r(x)=r(x)^{-1}Dr(x)\) when \(r(x)>0\). The real Hilbert-space gradient is the Riesz representative of the derivative; linearity of that identification gives \(\nabla\log r(x)=r(x)^{-1}\nabla r(x)\).



\noindent Lean declarations: \texttt{gradient\_log\_of\_pos}.

\section{gradient log add one of pos}


Adding a constant to a scalar function leaves its Fréchet derivative unchanged. Apply the gradient's derivative-based definition to obtain \(\nabla(\log r+1)=\nabla\log r\). The explicit differentiability/positivity arguments align the statement with the entropy-test calculus.



\noindent Lean declarations: \texttt{gradient\_log\_add\_one\_of\_pos}.

\section{inner gradient entropy test}


Rearrange the positive-density score identity to \(\nabla r=r\nabla\log r\). Pair this with \(\nabla(\log r+1)=\nabla\log r\) and use \(\langle s,s\rangle=\|s\|^2\). The Dirichlet integrand is therefore \(\langle\nabla r,\nabla(\log r+1)\rangle=r\|\nabla\log r\|^2\).



\noindent Lean declarations: \texttt{inner\_gradient\_entropy\_test}.

\section{entropy production of weak identity}


The strictly positive \(C^1\) density makes \(\psi=\log r+1\) a \(C^1\) test. The entropy derivative integrability is supplied, and the pointwise gradient identity converts its Dirichlet integrability to the supplied weighted Fisher integrability. Substitute this admissible test into the explicit weak equation to obtain \(\int\dot r(\log r+1)d\nu=-\int r\|\nabla\log r\|^2d\nu\).



\noindent Lean declarations: \texttt{entropy\_production\_of\_weak\_identity}.

\section{Test the weak equation with a noncompact function}


Use the same \(C^1\) entropy test in the forced weak equation. Rewrite the Dirichlet term by \texttt{inner\_gradient\_entropy\_test} and the forcing gradient by \texttt{gradient\_log\_add\_one\_of\_pos}. The three explicit integrability conditions justify the substitution, giving entropy production \(-I+\int r\langle e,\nabla\log r\rangle d\nu\). Admissibility of this noncompact test remains part of the weak-equation premise.



\noindent Lean declarations: \texttt{entropy\_production\_of\_forced\_weak\_identity}.

\end{document}
